% The Slide Definitions
%document
\documentclass[10pt]{beamer}
%theme
\usetheme{metropolis}
% packages
\usepackage{color}
\usepackage{listings}
\usepackage[ngerman]{babel}
\usepackage[utf8]{inputenc}
\usepackage{multicol}


% color definitions
\definecolor{mygreen}{rgb}{0,0.6,0}
\definecolor{mygray}{rgb}{0.5,0.5,0.5}
\definecolor{mymauve}{rgb}{0.58,0,0.82}

\lstset{
    backgroundcolor=\color{white},
    % choose the background color;
    % you must add \usepackage{color} or \usepackage{xcolor}
    basicstyle=\footnotesize\ttfamily,
    % the size of the fonts that are used for the code
    breakatwhitespace=false,
    % sets if automatic breaks should only happen at whitespace
    breaklines=true,                 % sets automatic line breaking
    captionpos=b,                    % sets the caption-position to bottom
    commentstyle=\color{mygreen},    % comment style
    % deletekeywords={...},
    % if you want to delete keywords from the given language
    extendedchars=true,
    % lets you use non-ASCII characters;
    % for 8-bits encodings only, does not work with UTF-8
    frame=single,                    % adds a frame around the code
    keepspaces=true,
    % keeps spaces in text,
    % useful for keeping indentation of code
    % (possibly needs columns=flexible)
    keywordstyle=\color{blue},       % keyword style
    % morekeywords={*,...},
    % if you want to add more keywords to the set
    numbers=left,
    % where to put the line-numbers; possible values are (none, left, right)
    numbersep=5pt,
    % how far the line-numbers are from the code
    numberstyle=\tiny\color{mygray},
    % the style that is used for the line-numbers
    rulecolor=\color{black},
    % if not set, the frame-color may be changed on line-breaks
    % within not-black text (e.g. comments (green here))
    stepnumber=1,
    % the step between two line-numbers.
    % If it's 1, each line will be numbered
    stringstyle=\color{mymauve},     % string literal style
    tabsize=4,                       % sets default tabsize to 4 spaces
    % show the filename of files included with \lstinputlisting;
    % also try caption instead of title
    language = [Sharp]C,
	showspaces = false,
	showtabs = false,
	showstringspaces = false,
	escapechar = ,
}

\def\ContinueLineNumber{\lstset{firstnumber=last}}
\def\StartLineAt#1{\lstset{firstnumber=#1}}
\let\numberLineAt\StartLineAt



\newcommand{\codeline}[1]{
	\alert{\texttt{#1}}
}


% Author and Course information
% This Document contains the information about this course.

% Authors of the slides
\author{\href{https://github.com/evilham}{@evilham}\\
\footnotesize{Basierend auf Folien von \href{https://github.com/satkowski}{@satkowski} (November 2016)}}

% Name of the Course
\institute{C\texttt{\#} Kurs}

% Fancy Logo
\titlegraphic{\hfill\includegraphics[height=1.25cm]{../templates/fsr_logo_cropped}}


% Presentation title
\title{Methoden (fortgeschritten)}
\date{\today}


\begin{document}

\maketitle

\begin{frame}{Gliederung}
	\setbeamertemplate{section in toc}[sections numbered]
	\tableofcontents
\end{frame}

% ----------------------- Extensions ------------------------------
\section{Extensions}
\begin{frame}{Extensions}
	\begin{itemize}
		\item Ermöglichen es vorhandenen Typen eigene Methoden hinzuzufügen
		\item Ohne eine abgeleitete Klasse zu erstellen oder den originalen Typen neu zu erstellen
		\item Sind statische Methoden, welche aber als Instanzmethoden auf den erweiterten Typ agieren
		\item Befinden sich selbst in einer extra statischen Klasse
		\item Syntax:
		\begin{itemize}
			\item \texttt{public static \alert{<Rückgabetyp> <Name>}( this \alert{<Erweiterter Typ> <Variablenname>}, \alert{[Übergabeparameter]} )\\ \{ \alert{<Code>} \}}
		\end{itemize}
	\end{itemize}
\end{frame}

\begin{frame}{Extensions}
	\lstinputlisting{resources/09_methoden_fortgeschritten_2/extensions.cs}
\end{frame}

% ----------------------- yield ------------------------------
\section{yield}
\begin{frame}{yield}
	\begin{itemize}
		\item Schlüsselwort
		\item Macht aus einer Methode, Operator oder get-Accessor einen Iterator
		\begin{itemize}
			\item Der Rückgabetyp muss \alert{IEnumerable}, \alert{IEnumerable<T>}, \alert{IEnumerator} oder \alert{IEnumerator<T>} sein
		\end{itemize}
		\item Syntax:
		\begin{itemize}
			\item \texttt{\alert{yield} return \alert{<Rückgabewert>};}
			\item \texttt{\alert{yield} break;}
		\end{itemize}
	\end{itemize}
\end{frame}

\begin{frame}{yield}
	\lstinputlisting{resources/09_methoden_fortgeschritten_2/yield.cs}	
\end{frame}

% ----------------------- Schlüsselwörter ------------------------------
\section{Schlüsselwörter}
\begin{frame}{Schlüsselwörter}
	\textbf{ref}\\
	\begin{itemize}
		\item Ermöglicht es einen Werttyp (wie z.B. int) als Referenztyp in eine Methode zu übergeben
		\item Wird dabei vor den Datentyp eines Übergabeparameters
		\item Und beim Aufruf vor die Variable geschrieben
	\end{itemize}
	\textbf{params}\\
	\begin{itemize}
		\item Ermöglicht es eine beliebige Anzahl von Übergabeparametern eines Typs einer Methode zu geben
		\item Diese werden in der Methode dann als Array genutzt
		\item Die Übergabe kann über eine mit Komma getrennte Liste oder über ein Array passieren
	\end{itemize}
\end{frame}

\begin{frame}{Schlüsselwörter}
	\lstinputlisting{resources/09_methoden_fortgeschritten_2/ref_param_creation.cs}
	\lstinputlisting{resources/09_methoden_fortgeschritten_2/ref_param_using.cs}
\end{frame}

\begin{frame}{Schlüsselwörter}
	\textbf{out}\\
	\begin{itemize}
		\item Eine Variable wird, wie bei \alert{ref}, als Verweis übergeben
		\item Jedoch muss diese Variable nun nicht deklariert sein
		\item Die Variable enthält zum Ende der Methode den zuletzt zugewiesenen Wert, es muss innerhalb der Methode mindestens eine Zuweisung erfolgen
		\item kann verwendet werden, wenn eine Methode mehr als einen Wert zurückgeben soll
	\end{itemize}
\end{frame}

\begin{frame}{Schlüsselwörter}
	\lstinputlisting{resources/09_methoden_fortgeschritten_2/out_creation.cs}
	\lstinputlisting{resources/09_methoden_fortgeschritten_2/out_using.cs}
\end{frame}

\end{document}
